\section{Results}
\label{sec:results}

\subsection{Historical transcript-input probe}
\label{sec:probe-results}

Table~\ref{tab:probe} reports an earlier controlled probe
(\S\ref{sec:experiments}) on
Greek$\rightarrow$English TEDx segments with English slides (M3 regime, human
references, $n{=}40$; Qwen3-32B, Local Agreement policy). The \emph{hard
stratum} contains segments where the audio-only baseline misses most reference
terminology (term recall $\leq 0.4$; $n{=}22$; the hard rate, $\sim$24\%, is
stable across 32B and 14B models --- a data property).
This experiment incrementally reveals a clean transcript rather than speech and
injects extracted terms as text. It measures whether advance textual hints can
help translation, but cannot establish acoustic disambiguation, raw visual
necessity, or the target persistent-evidence architecture.

\begin{table}[t]
\centering
\small
\begin{tabular}{lcccc}
\toprule
Condition & chrF & $\Delta$ pooled & $p$ & $\Delta$ hard \\
\midrule
A\; audio-only                & 38.8 & ---    & ---   & ---     \\
B\; slide terms (OCR)         & 35.2 & $-2.9$ & n.s.  & $+2.3$  \\
B\; slide terms (VLM)         & 39.9 & $-0.6$ & n.s.  & $+14.1$ \\
B\; VLM + need-gate (oracle)  & \textbf{45.3} & $\mathbf{+7.3}$ & $<.001$ & $+13.3$ \\
C\; oracle future terms       & 46.6 & $+7.0$ & 0.059 & $+16.3$ \\
D\; wrong-slide terms         & 38.6 & $-0.2$ & n.s.  & ---     \\
\bottomrule
\end{tabular}
\caption{Historical text-input diagnostic (el$\rightarrow$en, $n{=}40$;
paired bootstrap vs.\ A; vLLM serving). It compares textual hints and is not a
raw-vision result.}
\label{tab:probe}
\end{table}

\paragraph{Advance textual context has value on selected failures.}
Injecting the reference's own rare terms (C) --- a perfect-anticipation oracle
--- yields $+7.0$ chrF pooled and $+18.8$ on the hard stratum ($p{<}10^{-4}$),
with hard-stratum term recall rising $0.12\!\rightarrow\!0.41$ ($3.4\times$).
Gains concentrate on segments where the baseline fails on terminology and are
neutral elsewhere. The effect replicates across model scales (7B, 14B, 32B)
and across all three benchmark regimes (below).

\paragraph{Extraction quality matters within this text-hint protocol.} OCR
slide terms are neutral-to-harmful (480p fragments); VLM-read terms recover
$+14.1$ chrF on the selected hard segments. This establishes that a VLM can
produce better textual hints than the tested OCR pipeline, not that information
beyond strong OCR or structured text is necessary.

\subsection{Injection form: a three-step reversal}
\label{sec:injection-form}

Our initial runs suggested prompt hints disturb easy segments (motivating
need-gates and decode-level soft injection); the investigation closed in a
reversal that reshapes the method:

\begin{enumerate}[leftmargin=*]
\item \textbf{Apparent easy-segment harm} under vLLM serving affected
  \emph{all} hint conditions (including wrong-slide) with truncation
  signatures, and an oracle need-gate recovered $+7.3$ pooled ($p{<}.001$).
\item \textbf{Decode-level soft injection is benign but powerless}: uniform
  logit bias saturates at $+1.3$ hard ($+8$ bias), and trie-constrained
  shallow fusion --- boosting a term's continuation only along its prefix ---
  reaches only $+0.7$ pooled (termR $0.29$ vs.\ prompt's $0.46$).
  Terminology failures are \emph{phrase-planning} failures, not single-token
  near-misses; posterior logit nudges are too late and too local, while prompt
  hints act at the planning level.
\item \textbf{The root cause was serving nondeterminism}: batched vLLM serving
  returns different outputs for identical prefixes, stalling the
  longest-common-prefix commit. Under deterministic decoding the baseline
  jumps ($38.8\!\rightarrow\!62.7$ corpus chrF on the same segments) and
  prompt injection is simultaneously \emph{powerful and benign}: on 50
  14B-hard segments, $+6.5$ chrF ($p{<}0.001$), term recall
  $0.27\!\rightarrow\!0.46$, easy-$\Delta$ $\approx 0$.
\end{enumerate}

\noindent\textbf{Streaming-infrastructure verdict.} Deterministic (or drift-tolerant)
commit + prompt-level hint injection; need-gating demotes from a necessity to
a token-budget optimization; trie biasing is recorded as an informative
negative result. Practically: simultaneous pipelines built on serving-grade
inference must either pin deterministic kernels or use drift-tolerant commit
policies --- a deployment finding independent of the visual channel.

\subsection{Cross-regime mechanism evidence (H2)}
\label{sec:mechanism}

The En$\rightarrow$Zh control stratum (ACL 60/60, 5 talks, 60 segments,
official tagged terminology; slides are \emph{source}-language $\Rightarrow$
M1+M2 only) delivers the first direct test of the mechanism decomposition:

\begin{table}[t]
\centering
\small
\begin{tabular}{lcccc}
\toprule
Condition (En$\rightarrow$Zh) & chrF & $\Delta$ & termAcc & $\Delta$ hard \\
\midrule
none              & 26.3 & ---    & 0.42 & ---    \\
\textbf{slide (VLM, ungated)} & \textbf{28.3} & $+3.1$ & 0.44 & $+9.1$ \\
oracle terms      & 25.9 & $-1.1$ & 0.39 & $+3.8$ \\
wrong slide       & 23.1 & $-3.2$ & 0.32 & $+3.3$ \\
\bottomrule
\end{tabular}
\caption{S3 control (ACL 60/60, $n{=}60$, official tagged terms; vLLM-era
protocol, zh char-level LCP --- truncation caveat applies).}
\label{tab:s3}
\end{table}

\begin{itemize}[leftmargin=*]
\item \textbf{Copy-rate 0.00} (2/462): English slide strings never enter the
  Chinese output verbatim, versus the active copy channel in X$\rightarrow$En
  --- exactly H2's prediction that M3 (target-form supply) exists only when
  slide language equals the target language, while M2 (recognition support)
  gains persist ($+9.1$ hard) without copying.
\item \textbf{Real slides transfer ungated} on this stratum: ACL screen-recorded
  slides are dense (95\% of frames term-dense, $8\times$ mTEDx) and
  speech-aligned (speakers read their slides), so injection matches the
  model's natural phrasing.
\item \textbf{Faithfulness is direction-sensitive}: wrong-slide evidence is
  benign for 32B X$\rightarrow$En but harmful En$\rightarrow$Zh ($-3.2$);
  faithfulness metrics must be reported per direction.
\end{itemize}

\subsection{Learned need prediction (as budget optimization)}
Surface features of the source predict terminology failure at chance
(talk-held-out AUC 0.49--0.52); draft-translation logprob features reach AUC
0.62 --- the model's own uncertainty carries the signal --- preserving the full
hard-stratum gain end-to-end with imperfect precision (0.16). Under the
deterministic-decoding verdict this predictor is no longer required for
safety, but remains useful to prioritize hint budget in long sessions.

\subsection{Speech-input image control: content specificity is not established}

The first speech-input probe uses Qwen3-Omni-30B-A3B-Instruct on 206
Chinese-to-English lecture segments. It compares audio-only, the current slide
image, and a different slide image from the same lecture. Each segment is run
independently, so the visual encoder is repeated rather than cached.

\begin{table}[t]
\centering
\small
\begin{tabular}{lcc}
\toprule
Condition ($n{=}206$) & chrF & AL (audio chunks) \\
\midrule
audio only & 78.3 & 2.63 \\
current slide image & 80.7 & 2.42 \\
same-talk wrong slide image & 80.6 & 2.43 \\
\bottomrule
\end{tabular}
\caption{Speech-input diagnostic. Current versus same-talk wrong slide is
$+0.34$ chrF ($p{=}0.22$) and $-0.002$ AL chunks ($p{=}0.47$); the observed
gain does not establish use of local slide content.}
\label{tab:speech-image-control}
\end{table}

Both image conditions improve over audio-only, but current versus wrong is not
significant. The bounded conclusion is therefore negative: under this clean,
single-talk, per-segment protocol, the model does not demonstrate
segment-specific slide-content use. Same-talk domain priming, generic image
presence, and machine-reference bias remain possible. The persistent-evidence
study must add cross-talk, strong-OCR, text-equivalent, stale-time, and
once-per-slide cache controls before interpreting an image gain.

\subsection{Full-system results}
\todo{Main table for the persistent-evidence protocol: correct/stale/cross-talk
content controls; strong OCR/structured VLM/text-equivalent representations;
quality $\times$ computation-aware latency; cold, amortized, and on-path cost;
term/entity accuracy and visible-but-unspoken faithfulness.}
